\documentclass[10pt,a4paper]{amsart}
\usepackage[margin=19mm]{geometry}
\usepackage{amsmath,amssymb,mathtools}
\usepackage[scr=rsfs,cal=euler]{mathalfa}
\usepackage{xfrac}
\usepackage[unicode,hidelinks,bookmarks=false]{hyperref}
\newcommand{\ie}{i.e.\ }
\newcommand{\cf}{cf.\ }
\newcommand{\resp}{respectively\ }
\newcommand{\Subsection}{\subsection}
\providecommand{\nobreakdash}{\nobreak\hbox{-}}
\setlength{\emergencystretch}{2em}
\multlinegap0pt
\usepackage{amssymb}
\usepackage{tikz}
\newtheorem{theorem}{Theorem}
\newcommand{\C}{\mathbb{C}}
\def\Ind{\mathrm{Ind}}
\def\Spec{\mathrm{Spec}}
\newcommand{\Z}{\mathbb{Z}}
\title{A twisted Hecke algebra, then and now, and a Klein bottle of tempered representations}
\author{Anne-Marie Aubert, Roger Plymen}
\date{Source: arXiv:2603.03027v1 (CC BY 4.0)}
\begin{document}
\maketitle

\noindent\emph{Source and licence.} A twisted Hecke algebra, then and now, and a Klein bottle of tempered representations. Version arXiv:2603.03027v1 (CC BY 4.0), distributed by arXiv under the Creative Commons Attribution 4.0 International licence, http://creativecommons.org/licenses/by/4.0/, as recorded in the arXiv licence metadata for this exact version and shown on the abstract page https://arxiv.org/abs/2603.03027v1. Full licence notice: \texttt{licenses/arxiv-2603.03027-CC-BY-4.0.txt}.\par\smallskip

\noindent\textit{Editorial scope.} Counted unit: the article's theorem thm:simples-expanded (source0.tex lines 1047--1070). The article's complete original proof of it is delivered with it (source0.tex lines 1072--1201, one \texttt{proof} environment of 130 lines). No mathematics is written, repaired or re-derived: the statement and the proof are the article's own lines, in the article's own notation, and no retained line is substituted or re-flowed.  No further source-local context is needed: every cross-reference of every retained line resolves inside the two delivered spans.  Nothing was omitted from any retained span, so the package carries no omission record.  No retained line contains a citation, so the wrapper carries no bibliography and no bibliography entry.  No retained line contains a figure environment, an image-inclusion command or drawing code, so no image is delivered and the package declares no asset dependency.  Editorial formatting only, in addition: an AMS \texttt{amsart} preamble that restates the article's own declarations and macros used by the retained lines, namely the environments \texttt{theorem} on separate counters, together with the standard packages those lines need.  The retained environments print in this document as Theorem 1 in order of appearance, whereas the article prints the same items with the numbering of its own full document.  The complete native source of the article as distributed in the arXiv e-print for this exact version is bundled unchanged as \texttt{sources/195649/source0.tex}.
\begin{theorem}\label{thm:simples-expanded}
We have, for the twisted AFO-algebra,
\[
\C[\Z\rtimes\Z/2 \times \Z, \mu]
\;\cong\;
\C[s,X,X^{-1},Y,Y^{-1}]
\Big/\bigl(s^2-1,\; sY+Ys,\; sX-X^{-1}s,\; XY-YX\bigr),
\]
and the simple modules are the $2$--dimensional modules
\[
M_{w,z}\cong \C^2,\qquad
Y\mapsto\begin{pmatrix} w&0\\0&-w\end{pmatrix},\quad
s\mapsto\begin{pmatrix}0&1\\1&0\end{pmatrix},\quad
X\mapsto\begin{pmatrix} z&0\\0&z^{-1}\end{pmatrix},
\]
with parameters $(w,z)\in(\C^\times)^2$, satisfying
\[
M_{w,z}\;\cong\; M_{-w,z^{-1}},
\]
and every simple module is isomorphic to one of these.
%Consequently the parameter space for simple modules is \[(\C^\times\times\C^\times)\big/\bigl((w,z)\sim(-w,z^{-1})\bigr).\]
The primitive spectrum of $\C[\Gamma, \mu^{\mathcal{T}}]$  is naturally identified with $(\C^\times \times \C^\times) / (w,z) \sim (-w, z^{-1})$.   
In particular, the primitive spectrum is irreducible.
\end{theorem}

\begin{proof}
\emph{Presentation of the twisted group algebra.}
Write
\[
\Gamma=(\Z_A\rtimes \Z/2)\times \Z_B,
\]
where the generator $s$ of $\Z/2$ acts on $\Z_A$ by inversion.
Let $X$ denote the group-like element corresponding to $1\in \Z_A$
and $Y$ the group-like element corresponding to $1\in \Z_B$.
Then in the untwisted group algebra one has
\[
sXs^{-1}=X^{-1},\qquad XY=YX,\qquad s^2=1.
\]
The cocycle $\mu$ is trivial on $\Z_A\rtimes \Z/2$ and on $\Z_B$,
and is nontrivial only on the subgroup $\Z/2\times \Z_B$; explicitly,
with additive notation $(\varepsilon,m)\in \Z/2\times\Z_B$,
\[
\mu\bigl((\varepsilon,m),(\delta,n)\bigr)=(-1)^{\varepsilon n}.
\]
In particular, if we write $N_g$ for the basis element attached to $g\in\Gamma$,
then in $\C[\Gamma,\mu]$ we have
\[
N_sN_Y=\mu(s,Y)N_{sY}=(-1)N_{sY},\qquad
N_YN_s=\mu(Y,s)N_{Ys}=(+1)N_{Ys}.
\]
Since $Ys=sY$ in the underlying group, this yields the twisted relation
\[
N_sN_Y=-\,N_YN_s.
\]
Thus $\C[\Gamma,\mu]$ is generated by $s,X^{\pm1},Y^{\pm1}$ with relations
\[
s^2=1,\qquad sX=X^{-1}s,\qquad sY=-Ys,\qquad XY=YX,
\]
which gives the stated presentation.

\medskip
\emph{No one-dimensional modules.}
Let $M$ be a $1$--dimensional module. Since $Y$ is invertible in the algebra,
its image in $\mathrm{End}(M)\cong \C$ must be a nonzero scalar, say $Y\mapsto \lambda\in\C^\times$.
But the relation $sY=-Ys$ forces
\[
s\lambda = -\lambda s.
\]
In $\C$ this implies $2\lambda s=0$, hence $s=0$, contradicting $s^2=1$.
Therefore \emph{no} $1$--dimensional representations exist.

\medskip
\emph{Reduction to Clifford theory over a commutative subalgebra.}
Let
\[
B:=\C[X^{\pm1},Y^{\pm1}]\subset \C[\Gamma,\mu].
\]
Then $B$ is a commutative Laurent polynomial algebra in two variables.
Conjugation by $s$ defines an involutive automorphism $\alpha$ of $B$:
\[
\alpha(X)=X^{-1},\qquad \alpha(Y)=-Y.
\]
Equivalently, for any character $\chi:B\to\C$ determined by
\[
\chi(X)=z\in\C^\times,\qquad \chi(Y)=w\in\C^\times,
\]
one has
\[
(\chi\circ\alpha)(X)=z^{-1},\qquad (\chi\circ\alpha)(Y)=-w.
\]
Thus the $\Z/2$--action on $\Spec(B)\cong(\C^\times)^2$ is
\[
(w,z)\longmapsto(-w,z^{-1}).
\]
Note that this action has \emph{no fixed points} on $(\C^\times)^2$:
a fixed point would require $z=z^{-1}$ and $w=-w$, hence $w=0$, impossible.

\medskip
\emph{Construction of the $2$--dimensional simple modules.}
Fix $(w,z)\in(\C^\times)^2$ and let $\C_{w,z}$ be the $1$--dimensional $B$--module
with $X\mapsto z$ and $Y\mapsto w$.
Form the induced module
\[
M_{w,z}:=\Ind_{B}^{\C[\Gamma,\mu]}\C_{w,z}\;\cong\;\C[\Gamma,\mu]\otimes_B \C_{w,z}.
\]
Because the $\Z/2$--orbit of $(w,z)$ has size $2$, standard Clifford theory
(or a direct computation below) shows $M_{w,z}$ is $2$--dimensional and simple.

Concretely, let $v:=1\otimes 1\in M_{w,z}$ and set $v':=s\cdot v$.
Then $\{v,v'\}$ is a basis of $M_{w,z}$.
Using $s^2=1$ we have $s\cdot v=v'$ and $s\cdot v'=v$, so
\[
s\mapsto \begin{pmatrix}0&1\\[2pt]1&0\end{pmatrix}.
\]
Next, $Xv=zv$ by definition.
Also,
\[
Xv' = X(sv) = (Xs)v = (sX^{-1})v = s(z^{-1}v)=z^{-1}v',
\]
so
\[
X\mapsto \begin{pmatrix} z&0\\[2pt]0&z^{-1}\end{pmatrix}.
\]
Similarly, $Yv=wv$ and
\[
Yv' = Y(sv) = -(sY)v = -s(wv)= -wv',
\]
hence
\[
Y\mapsto \begin{pmatrix} w&0\\[2pt]0&-w\end{pmatrix}.
\]
These matrices satisfy the defining relations, so $M_{w,z}$ is indeed a representation.

\medskip
\emph{Isomorphism relation $M_{w,z}\cong M_{-w,z^{-1}}$.}
The $\Z/2$--action on characters of $B$ sends $(w,z)$ to $(-w,z^{-1})$.
Equivalently, conjugating the above matrices by the $s$--matrix swaps the two basis vectors
and sends $(w,z)$ to $(-w,z^{-1})$.
Thus $M_{w,z}\cong M_{-w,z^{-1}}$.

\medskip
\emph{Exhaustion: every simple module is one of these.}
Let $M$ be a simple $\C[\Gamma,\mu]$--module.
Restrict $M$ to the commutative subalgebra $B$.
Choose a maximal ideal $\frak m\subset B$ in the support of $M$; equivalently,
choose a character $\chi:B\to\C$ occurring in $M$, with $\chi(X)=z\neq 0$ and $\chi(Y)=w\neq 0$
(since $X,Y$ are invertible in the algebra).
Because the $\Z/2$--action on $\Spec(B)$ is free, the $\chi$--isotypic component generates a
$\C[\Gamma,\mu]$--submodule of dimension $2$ (spanned by $v$ and $sv$ as above), and simplicity
forces $M$ to be isomorphic to this induced module $M_{w,z}$.
Therefore every simple module is $2$--dimensional and isomorphic to some $M_{w,z}$, with the
sole identification $(w,z)\sim(-w,z^{-1})$.

This completes the proof.
\end{proof}

\end{document}
